\documentclass[11pt,a4paper]{article}
\usepackage[margin=25mm]{geometry}
\usepackage{amsmath,amssymb,amsthm,microtype}
\usepackage[colorlinks=true,citecolor=blue,urlcolor=blue]{hyperref}
\newtheorem{theorem}{Theorem}
\newtheorem{lemma}{Lemma}
\newtheorem{proposition}{Proposition}
\title{\textbf{Exact Low-Area Center-Flux Surfaces and Strong-Coupling Coefficients}\\in Four-Dimensional Lattice Gauge Theory}
\author{Tony Newton\\Newton Astro Labs\thanks{Newton Astro Labs is the trading name under which Tony Newton conducts independent computational research in the United Kingdom; it is not a limited company.}\\London, UK\\\texttt{tony.newton79@gmail.com}}
\date{}
\begin{document}
\maketitle
\begin{abstract}
We classify the first two nonempty mod-two closed plaquette-surface strata of the four-dimensional cubic lattice. Every finite closed surface with fewer than twelve faces is the six-face boundary of an elementary cube or the ten-face symmetric difference of two distinct elementary cubes sharing a face. The second presentation is unique. Rooted multiplicities are four at area six and sixty at area ten, with the latter split into twenty same-orientation and forty mixed-orientation cases. We derive the first two coefficients of the reduced finite-volume abelian lattice-gauge partition function, and explain the corresponding necessary parity rule in SU(2) character expansions. These are known-type strong-coupling results with independently reproduced exact counts, not a new continuum Yang--Mills mass-gap theorem.
\end{abstract}
\section{Definitions and prior results}
Let \(C_k(\mathbb Z^d;\mathbb F_2)\) denote finitely supported cubical chains. A two-chain \(S\) of unit plaquettes is closed when \(\partial S=0\), equivalently each lattice edge has even incidence. The area \(|S|\) counts distinct unoriented plaquettes. The closed-surface strong-coupling expansion is classical \cite{wegner,kogut,drouffe}; related rigorous four-dimensional small-surface results appear in \cite{forsstrom}. We do not claim novelty of the initial cube terms.

\section{Minimum cubical-cycle weight and parity}
\begin{lemma}\label{lem:min}
For \(1\le k<d\), any nonempty finitely supported mod-two cubical \(k\)-cycle in \(\mathbb Z^d\) contains at least \(2(k+1)\) cells. The minima for \(k=1,2,3\) are therefore \(4,6,8\).
\end{lemma}
\begin{proof}
A nonempty finite one-cycle in the bipartite nearest-neighbor grid graph contains a simple cycle of at least four edges. For induction, select any coordinate direction \(j\) occurring in the support of a \(k\)-cycle. The \(j\)-containing \(k\)-cells in the slab \(n\le x_j\le n+1\) project to a transverse \((k-1)\)-chain \(C_n\). The \(j\)-parallel part of the boundary condition implies \(\partial C_n=0\), so each nonzero \(C_n\) has at least \(2k\) cells by induction. Let \(H_n\) be the horizontal \(k\)-chain at coordinate \(x_j=n\). The other part of the boundary condition is
\[
\partial H_n=C_{n-1}+C_n\quad\text{over }\mathbb F_2.
\]
The first and last nonzero slabs each require a nonempty horizontal cap. Hence at least \(2k+2\) cells occur in total.
\end{proof}
\begin{lemma}\label{lem:even}
Every finite mod-two closed plaquette two-chain in \(\mathbb Z^d\) has even area.
\end{lemma}
\begin{proof}
For a positively oriented edge based at \(x\) along direction \(j\), define the mod-two one-cochain
\(\eta(x,j)=\sum_{i<j}x_i\bmod2\).
On every coordinate plaquette with axes \(i<j\), \(\delta\eta=1\): the \(j\)-edge contribution changes by one along \(i\) while the \(i\)-edge contribution does not change along \(j\). Thus for \(\partial S=0\),
\[
|S|\bmod2=\langle\delta\eta,S\rangle
=\langle\eta,\partial S\rangle=0.
\]
\end{proof}


\section{Complete classification below area twelve}
\begin{theorem}\label{thm:small}
For \(d\ge3\), every nonempty finite mod-two closed plaquette surface in \(\mathbb Z^d\) with fewer than twelve plaquettes is either (i) the six-face boundary of an elementary three-cube, or (ii) the ten-face symmetric difference of boundaries of two distinct elementary three-cubes sharing one face. The unordered pair in (ii) is unique.
\end{theorem}
\begin{proof}
Choose a coordinate direction \(j\) used by the surface. As in Lemma~\ref{lem:min}, the active \(j\)-slabs carry transverse finite one-cycles \(C_n\) with even length at least four, and the horizontal chains satisfy \(\partial H_n=C_{n-1}+C_n\). If \(m\) slabs are active, their combined faces plus the two end caps give at least \(4m+2\) plaquettes. Hence area at most ten permits only one or two active slabs.

For one slab with \(|C_n|=4\), its loop bounds a unique unit square \(Q\). Each cap \(H\) having \(\partial H=\partial Q\) equals \(Q\), of area one, or has area at least five: otherwise the nonzero two-cycle \(H\triangle Q\) would have fewer than six faces, contradicting Lemma~\ref{lem:min}. The only totals at most ten are \(4+1+1=6\), the boundary of one cube extruding \(Q\), or \(4+1+5=10\). In the latter case the larger cap differs from \(Q\) by a six-face closed cycle, necessarily another cube boundary, with its face \(Q\) cancelled.

For one slab with \(|C_n|=6\), the caps need at least two squares each. Equality at ten faces requires two two-square fillings of the loop. Such a filling consists of two unit squares sharing one edge. The two fillings coincide, because their symmetric difference would otherwise be a nonempty closed two-cycle with at most four faces. Extruding these squares along \(j\) gives two elementary cubes sharing the edge-times-interval plaquette. If \(|C_n|\ge8\), the two caps each require at least two faces, giving total area at least twelve.

With two active slabs, equality at \(4+4+1+1=10\) forces adjacent slabs with identical elementary square loops, single-square exterior caps, and no intermediate faces. This is a pair of consecutive cubes with their common face cancelled. Separated active slabs would require additional intermediate caps. The three-slab lower bound is fourteen. Lemma~\ref{lem:even} excludes eleven plaquettes.

If two different cube pairs produced the same ten-face surface, their difference as mod-two three-chains would be a nonzero closed three-chain supported on at most four cubes. This contradicts the eight-cell lower bound for three-cycles in Lemma~\ref{lem:min}. Conversely two distinct elementary cubes sharing one face have a closed symmetric-difference boundary with \(6+6-2=10\) faces.
\end{proof}

% MORE_CONTENT_HERE
